% \todo{Bu Turkce ozeti ana dilde gozden gecirip kesinlestirin (taslak ceviridir).}
\noindent
Bu tez, eser miktarda diazot monoksitin (\nto) tespiti için bir kuvars-geliştirilmiş fotoakustik spektroskopi (QEPAS) algılayıcısının tasarımını, kurulumunu ve deneysel karakterizasyonunu sunmaktadır. Algılayıcı, uyarma kaynağı olarak darbeli ve ayarlanabilir bir kuantum kademeli lazer (QCL) kullanır; lazer, bir akustik algılama modülü içine yerleştirilmiş bir kuvars diyapazonunun rezonans frekansında genlik modülasyonu modunda çalıştırılır ve fotoakustik sinyal bir kilitlemeli yükselteç ile geri kazanılır. Lazer, mevcut kaynağın ayar aralığına düşen, \SI{7.7}{\micro\meter} civarındaki \nto\ $\nu_1$ bandındaki bir çizgiye ayarlanır. Çalışma; optik, akustik ve elektronik alt sistemlerin entegrasyonunu ve hizalanmasını, kısa darbe genişliklerinde kilit kaybına çözüm içeren harici tetiklemeli bir genlik modülasyonu şemasını ve ışın iletim yolu için dalga boyuna özgü bir optik güç bütçesini kapsar. Kuvars diyapazon, kilitlemeli yükselteç zinciri ve lazer ayrı ayrı karakterize edilir; \nto\ tespit edilir ve konsantrasyona karşı kalibre edilerek algılayıcının başarımı duyarlılık, normalize edilmiş gürültü eşdeğeri soğurma, tespit edilebilir en düşük konsantrasyon ve Allan sapması analizinden elde edilen kararlılık açısından nicelenir. \todo{Ölçümler tamamlandığında nihai temel sayısal değerleri ekleyin.} Elde edilen cihaz, gelecekte eser miktarda uçucu organik bileşiklerin tespiti için doğrulanmış bir platform sağlar.
\\

\noindent
\textbf{\textit{Anahtar Kelimeler}}: kuvars-geliştirilmiş fotoakustik spektroskopi (QEPAS), diazot monoksit, darbeli kuantum kademeli lazer, genlik modülasyonu, eser gaz tespiti
